\documentclass[10pt]{article}
\usepackage[a4paper,margin=0.5in]{geometry}
\usepackage{amsmath,amssymb,amsfonts,mathtools}
\usepackage[dvipsnames]{xcolor}
\usepackage{array,booktabs,longtable,float}
\usepackage[colorlinks=true,linkcolor=blue]{hyperref}
\setlength{\emergencystretch}{3em}
\newcommand{\KK}{\mathcal K}
\title{Transverse divergences of the terms in your \texttt{Other\_terms.tex}}
\date{}
\begin{document}
\maketitle
\noindent Your file without Group I Row II part 1 and Group II Row IV (both inside \verb|\begin{comment}|). Question: does any term diverge when transverse coordinates come together (UV), or anywhere else in the transverse plane? $\KK(X)=X/X^2$, $a=\frac{k^+}{\vee}$, $U^{bc}=U^{cb}$.

\section{Answer}
\begin{center}\fbox{\begin{minipage}{0.96\textwidth}
\textbf{Flagged.}
\begin{enumerate}
\item[\textbf{F1}] \textbf{Group I Row V, $d^3N_{2b\,\rm UV}$: UV divergent.} The $\frac{2}{\epsilon}$ is the dimensionally regularized UV divergence of 2b at $x\to z$ (the two gluons of the loop at one point). After it, $\log\frac{(y-z)^2\mu^2}{4}$ times $\frac{(y-z)^i}{(y-z)^2}[U(y)-U(z)]$ is integrable at $y\to z$. In this file its partner is gone: the $\frac{11}{3}$ term of Group I Row II part 1, which cancels its $\beta_0$ pole, is commented out.
\item[\textbf{F2}] \textbf{Group I Row III: each of the parts 2, 3a, 3b is divergent; only the sum of the four parts is finite.}
\begin{itemize}
\item At $z\to w$, each of parts 2, 3a, 3b grows like $r^{-2}$ (logarithmic divergence).
\item Across the curve $(x-z)^2=(x-w)^2$ (it passes through $z=w$), the separate lines go like $[(x-z)^2-(x-w)^2]^{-2}$, which is not integrable.
\item Across the segment where $w$ lies between $x$ and $z$, the $\mathcal A_0$ lines go like $1/{\rm distance}$, which is not integrable either.
\item The sum of parts 1, 2, 3a, 3b is regular on both curves, and its worst UV behaviour is $r^{-1}$ (at $x\to w$), which is integrable.
\end{itemize}
Row III must be kept as one expression.
\item[\textbf{F3}] \textbf{Group I Row V, $d^3N_{2b\,\rm finite}$: separate lines divergent at $x\to z$, the sum finite.} The $d_\perp$ line, the line without the plus prescription, and the $\frac{11}{6}$ line each go like $(x-z)^{-2}$. The $\frac{11}{6}$ line is the subtraction that removes this divergence.
\item[\textbf{F4}] \textbf{Group II Row I, $d^3N_{2f}$: each of its three lines is divergent, the sum finite.} Line 2, $\big[\frac1\vee-\frac1{k^+}\big]\frac{1}{(z-w)^2(z-w')^2}$, diverges logarithmically at $z\to w$ and at $z\to w'$. Line 1 diverges at $z\to w$ and line 3 at $z\to w'$, and they cancel it. Keep the three lines together.
\item[\textbf{F5}] \textbf{Group I Row II, $d^3N_{2b}$ (momentum form, $R^i$): UV finite, but infrared divergent} (not UV). Two regions are affected.
\begin{itemize}
\item At $p\to0$: $R^i\to-2k^i\log\frac{p^2}{k^2}$.
\item At $q=k-p\to0$: $R^i\to\frac{2k^2k^i}{q^2}\log\frac{k^2}{q^2}$.
\end{itemize}
With the $\frac{1}{p^2k^2}$, $\int d^2p$ diverges like $\log^2$ at both points. These are the regions $|x-y|\to\infty$. In addition, $\int d^2k$ uses the same letter $k$ as the measured gluon.
\end{enumerate}
\textbf{Not divergent} (UV and on curves): every other equation, i.e.\ Group I Rows III (sum) and IV, Row V first part and $d^3N_{2a}$, Group II Row I (every equation; separate lines of the second part need a principal value across the curve $(x'-w')^2(x-z)^2=(x-w)^2(x'-z)^2$), Group III Rows I, II, III, V, VI.
\end{minipage}}\end{center}

\section{Method}
\begin{itemize}
\item Your equations are evaluated pointwise as written (each equation split into its lines, \texttt{terms.py}). Wilson lines and charges are smooth functions of position:
\begin{itemize}
\item $U(p)$ is the adjoint of $V(p)=W(p)\,{\rm diag}(1,-1,-1)\,W(p)^\dagger$, with $W(p)=e^{i\theta_a(p)t^a}$; this is an SU(3) matrix and $U=U^T$.
\item $\theta_a(p)$ and $\rho^a(p)$ are sums of sines with wave numbers of order 1.
\end{itemize}
A colour factor that vanishes when two points coincide therefore vanishes linearly, as it does in the real integral.
\item \textbf{UV.} For every pair of points $P,Q$, set $P=Q+r\,\hat n(\varphi)$ and average the integrand over $\varphi$, giving $A(r)$ at $r=10^{-1.5}\dots10^{-4.5}$. The exponent $s$ in $A\sim r^s$ decides: $\int d^2r$ diverges iff $s\le-2$.
\begin{itemize}
\item Where a singular curve crosses the circle, the $\varphi$ average is a principal value.
\item Every equation was scanned, and every row with several equations as a whole.
\item Two field and point configurations were used. Terms with $p^+$ or $q^+$ integrals are evaluated at fixed $p^+$ ($0.37k^+$; the plus prescriptions subtract at $p^+=0$), on the pairs of points that can diverge (Sec.~\ref{sec:pc}).
\item In Group I Row V $d^3N_{2b\,\rm finite}$ the lines have different longitudinal variables ($p^+$, $q^+$, none), so they can cancel only after the integrals. There the $p^+$ and $q^+$ integrals are done (16-point Gauss--Legendre on $[0,k^+]$, with the orientation and $\frac{1}{2\pi}$ of each), \texttt{intcheck.py}.
\end{itemize}
\item \textbf{Singular curves.} A point is moved across the curve; from $g(\pm\delta)$ take $E=\frac{g_++g_-}{2}\sim\delta^{e_E}$ and $O=\frac{g_+-g_-}{2}\sim\delta^{e_O}$.
\begin{itemize}
\item Integrable: $e_E>-1$ and $e_O>-1$.
\item Principal value: $e_O=-1$.
\item Divergent: $e_E\le-1$ or $e_O\le-2$.
\end{itemize}
\end{itemize}

\section{UV: coincident points}
$s$ = exponent of $A(r)$ at the worst pair (all configurations and $p^+$ values); $s\le-2$ = divergent.
{\small
\begin{longtable}{p{5.2cm}cp{1.6cm}p{2.4cm}p{1.4cm}p{4.2cm}}
\toprule
equation & worst $s$ & pair & divergent pairs & verdict & scan\\
\midrule
Group I, Row III, eq.~1 & -0.12 & $x\to w$ & --- & finite & {\footnotesize all pairs, 2 configurations}\\
Group I, Row III, eq.~2 & -2.07 & $z\to w$ & $z\to w$ & \textcolor{red}{divergent} & {\footnotesize all pairs, 2 configurations}\\
Group I, Row III, eq.~3 & -2.07 & $z\to w$ & $z\to w$ & \textcolor{red}{divergent} & {\footnotesize all pairs, 2 configurations}\\
Group I, Row III, eq.~4 & -2.00 & $z\to w$ & $z\to w$ & \textcolor{red}{divergent} & {\footnotesize all pairs, 2 configurations}\\
Group I, Row III \textbf{row total} & -1.01 & $x\to w$ & --- & finite & {\footnotesize all pairs, 2 configurations}\\
Group I, Row IV, eq.~1 & 0.00 & $x\to y$ & --- & finite & {\footnotesize all pairs, 2 configurations}\\
Group I, Row V, 1st part, eq.~1 & -0.00 & $x\to y$ & --- & finite & {\footnotesize selected pairs, $p^+=0.37k^+$, line by line}\\
Group I, Row V, 2nd part, eq.~1 & -0.00 & $x\to y$ & --- & finite & {\footnotesize selected pairs, $p^+=0.37k^+$, line by line}\\
Group I, Row V, 2nd part, eq.~3 & -0.00 & $x\to z$ & --- & finite & {\footnotesize p$^+$, q$^+$ integrals done; lines 1, 3, 4 diverge at $x\to z$}\\
Group II, Row I, 1st part, eq.~1 & -0.00 & $x\to y$ & --- & finite & {\footnotesize all pairs, 2 configurations}\\
Group II, Row I, 1st part, eq.~2 & -0.00 & $x\to y$ & --- & finite & {\footnotesize all pairs, 2 configurations}\\
Group II, Row I, 2nd part, eq.~1 & -0.12 & $x\to w$ & --- & finite & {\footnotesize all pairs, 2 configurations}\\
Group II, Row I, 2nd part, eq.~2 & -0.13 & $x'\to w'$ & --- & finite & {\footnotesize all pairs, 2 configurations}\\
Group II, Row I, 2nd part, eq.~3 & -0.00 & $x\to z$ & --- & finite & {\footnotesize all pairs, 2 configurations}\\
Group II, Row I, 2nd part, eq.~4 & -0.13 & $x'\to w'$ & --- & finite & {\footnotesize all pairs, 2 configurations}\\
Group II, Row I, 2nd part, eq.~5 & -0.00 & $x\to z$ & --- & finite & {\footnotesize all pairs, 2 configurations}\\
Group II, Row I, 2nd part, eq.~6 & -0.01 & $z\to w'$ & --- & finite & {\footnotesize all pairs, 2 configurations; line 1: $z\to w$; line 2: $z\to w$, $z\to w'$; line 3: $z\to w'$}\\
Group II, Row I \textbf{row total} & -0.00 & $x\to y$ & --- & finite & {\footnotesize all pairs, 2 configurations}\\
Group III, Row I, eq.~1 & 0.00 & $x\to y$ & --- & finite & {\footnotesize all pairs, 2 configurations}\\
Group III, Row II, eq.~1 & 0.00 & $x\to y$ & --- & finite & {\footnotesize all pairs, 2 configurations}\\
Group III, Row III, 1st part, eq.~1 & -0.00 & $x\to z$ & --- & finite & {\footnotesize selected pairs, $p^+=0.37k^+$, line by line}\\
Group III, Row III, 2nd part, eq.~1 & 0.00 & $x\to z$ & --- & finite & {\footnotesize selected pairs, $p^+=0.37k^+$, line by line}\\
Group III, Row V, eq.~1 & -0.00 & $x\to z$ & --- & finite & {\footnotesize selected pairs, $p^+=0.37k^+$, line by line}\\
Group III, Row VI, eq.~1 & -0.00 & $x\to z$ & --- & finite & {\footnotesize selected pairs, $p^+=0.37k^+$, line by line}\\
\bottomrule
\end{longtable}}
Exponents near $0$ or $-1$ are bounded or $1/r$ singularities, both integrable. Group I Row II $d^3N_{2b}$ (momentum integrals) and Row V $d^3N_{2b\,\rm UV}$ (explicit pole) are not scanned; see F5 and F1.

\section{Singular curves}
{\small
\begin{longtable}{p{4.3cm}p{3.6cm}p{3.4cm}p{3cm}p{2.2cm}}
\toprule
equation & curve & separate lines & sum: $e_E$, $e_O$ & sum\\
\midrule
Group I, Row III, eq.~1 & $(x-z)^2=(x-w)^2$ & \textcolor{red}{DIVERGENT}, integrable & $e_E=0.00$, $e_O=-1.00$ & principal value\\
Group I, Row III, eq.~2 & $(x-z)^2=(x-w)^2$ & \textcolor{red}{DIVERGENT}, principal value & $e_E=-2.11$, $e_O=-1.17$ & \textcolor{red}{DIVERGENT}\\
Group I, Row III, eq.~3 & $(x-z)^2=(x-w)^2$ & \textcolor{red}{DIVERGENT}, principal value & $e_E=-2.11$, $e_O=-1.10$ & \textcolor{red}{DIVERGENT}\\
Group I, Row III, eq.~3 & $w$ on the segment $[x,z]$ ($\mathcal A_0$) & \textcolor{red}{DIVERGENT} & $e_E=-1.00$, $e_O=-0.01$ & \textcolor{red}{DIVERGENT}\\
Group I, Row III, eq.~4 & $w$ on the segment $[x,z]$ ($\mathcal A_0$) & \textcolor{red}{DIVERGENT} & $e_E=-1.00$, $e_O=-0.00$ & \textcolor{red}{DIVERGENT}\\
Group II, Row I, 2nd part, eq.~1 & $(x-z)^2=(x-w)^2$ & integrable & $e_E=-0.00$, $e_O=1.00$ & integrable\\
Group II, Row I, 2nd part, eq.~1 & $(x'-w')^2(x-z)^2=(x-w)^2(x'-z)^2$ & integrable & $e_E=0.00$, $e_O=1.00$ & integrable\\
Group II, Row I, 2nd part, eq.~2 & $(x'-w')^2(x-z)^2=(x-w)^2(x'-z)^2$ & integrable, principal value & $e_E=-0.00$, $e_O=1.00$ & integrable\\
Group II, Row I, 2nd part, eq.~3 & $(x-z)^2=(x-w)^2$ & integrable & $e_E=0.00$, $e_O=1.00$ & integrable\\
Group II, Row I, 2nd part, eq.~3 & $(x'-w')^2(x-z)^2=(x-w)^2(x'-z)^2$ & principal value & $e_E=-0.00$, $e_O=1.00$ & integrable\\
Group II, Row I, 2nd part, eq.~4 & $(x-z)^2=(x-w)^2$ & integrable & $e_E=0.00$, $e_O=1.00$ & integrable\\
Group II, Row I, 2nd part, eq.~4 & $(x'-w')^2(x-z)^2=(x-w)^2(x'-z)^2$ & principal value & $e_E=-0.00$, $e_O=1.00$ & integrable\\
Group II, Row I, 2nd part, eq.~4 & $(x'-w')^2=(x'-z)^2$ & integrable & $e_E=-0.00$, $e_O=1.00$ & integrable\\
Group II, Row I, 2nd part, eq.~5 & $(x-z)^2=(x-w)^2$ & integrable & $e_E=-0.00$, $e_O=1.00$ & integrable\\
Group II, Row I, 2nd part, eq.~5 & $(x'-w')^2(x-z)^2=(x-w)^2(x'-z)^2$ & principal value & $e_E=0.00$, $e_O=1.00$ & integrable\\
Group II, Row I, 2nd part, eq.~5 & $(x'-w')^2=(x'-z)^2$ & integrable & $e_E=-0.00$, $e_O=1.00$ & integrable\\
Group II, Row I, 2nd part, eq.~6 & $(x-z)^2=(x-w)^2$ & integrable & $e_E=2.00$, $e_O=1.00$ & integrable\\
Group II, Row I, 2nd part, eq.~6 & $(x'-w')^2(x-z)^2=(x-w)^2(x'-z)^2$ & principal value & $e_E=-0.00$, $e_O=1.00$ & integrable\\
Group II, Row I, 2nd part, eq.~6 & $(x'-w')^2=(x'-z)^2$ & integrable & $e_E=1.80$, $e_O=1.00$ & integrable\\
Group III, Row III, 2nd part, eq.~1 & $(x-z)^2=(x-w)^2$ & integrable & $e_E=0.00$, $e_O=1.00$ & integrable\\
Group I, Row III \textbf{row total} & $(x-z)^2=(x-w)^2$ & \textcolor{red}{DIVERGENT}, integrable, principal value & $e_E=0.00$, $e_O=0.99$ & integrable\\
Group I, Row III \textbf{row total} & $w$ on the segment $[x,z]$ ($\mathcal A_0$) & \textcolor{red}{DIVERGENT} & $e_E=-0.00$, $e_O=-0.00$ & integrable\\
Group II, Row I \textbf{row total} & $(x-z)^2=(x-w)^2$ & integrable & $e_E=0.00$, $e_O=1.00$ & integrable\\
Group II, Row I \textbf{row total} & $(x'-w')^2(x-z)^2=(x-w)^2(x'-z)^2$ & integrable, principal value & $e_E=-0.00$, $e_O=1.00$ & integrable\\
Group II, Row I \textbf{row total} & $(x'-w')^2=(x'-z)^2$ & integrable & $e_E=-0.00$, $e_O=1.00$ & integrable\\
Group III, Row III \textbf{row total} & $(x-z)^2=(x-w)^2$ & integrable & $e_E=0.00$, $e_O=1.00$ & integrable\\
\bottomrule
\end{longtable}}

\section{Analytic checks}\label{sec:pc}
\paragraph{Power counting.} In two dimensions $\int d^2r\,r^{-n}$ converges at $r\to0$ for $n<2$. When $m$ points collapse together the measure is $r^{2m-3}dr$. Each $\KK(u-v)$ contributes $r^{-1}$ when $u\to v$. A colour bracket like $U(z)-U(y)$ adds a factor $r$ when $y\to z$.
\begin{itemize}
\item \textbf{Group III Rows III, V, VI, line 1} ($p^+$ integral) has $\KK^m(y'-z)\,\KK^{k'}(x'-y')\,\KK^i(y-w)\,\KK^j(x-z)$. Each kernel joins a different pair, so every two-point limit is $r^{-1}$. A collapse of $m$ points brings at most $m-1$ kernels, i.e.\ $r^{-(m-1)}$ against $r^{2m-3}$. The same holds for line 2 ($\KK(x'-z)$ in place of $\KK(x'-y')$) and line 3. Hence these rows are finite.
\item \textbf{Group III Row III, second part.} $p^+(x-z)^2+k^+(x-w)^2$ vanishes only if $x,z,w$ all coincide. Its $\frac{1}{(z-w)^2}$ always multiplies $(z-w)$ or $(x-z)^2-(x-w)^2$, which is $O(r)$ at $z\to w$.
\item \textbf{Group I Row V, $d^3N_{2a}$.} At $x\to z$ the $d_\perp$ line goes like $\frac{(z-x)^i}{(x-z)^2}\frac{(y-x)^2-(y-z)^2}{(x-z)^2}\sim r^{-2}$, with a non-zero angular average. But its colour factor vanishes there: $f^{adc}U^{db}(z)U^{ce}(z)=f^{hbe}U^{ah}(z)$ and $f^{fbe}f^{hbe}=N_c\delta^{fh}$. So
\begin{equation*}
\big[-f^{adc}f^{fbe}U^{db}(x)U^{ce}(z)+N_cU^{af}(x)\big]_{x=z}=-N_cU^{af}(z)+N_cU^{af}(z)=0,
\end{equation*}
and the integrand is $O(r^{-1})$. The same holds for the $q^+$ line.
\item \textbf{Group I Row V, $d^3N_{2b\,\rm finite}$.} Its colour factor $N_c[U(y)-U(x)]$ does not vanish at $x\to z$.
\begin{itemize}
\item The $d_\perp$ line and the $q^+$ line without the plus prescription go like $(x-z)^{-2}$, with a non-zero average.
\item In the $[\frac{k^+}{q^+}]_+$ line the leading term does not depend on $q^+$ at $x=z$, so the plus prescription removes it.
\item The $\frac{11}{6}$ line, with $\frac{(y-x)\cdot(y-z)}{(y-x)^2(x-z)^2}\to\frac{1}{(x-z)^2}$, cancels the rest. This is the subtraction of \texttt{rowV/Part2b\_MSbar.pdf}: the sum is finite for every $p^+$.
\end{itemize}
\end{itemize}

\paragraph{Group I Row III, the $\mathcal A_0$ segment.} $\mathcal A_0=\int_0^1\frac{dt}{|(x-w)-t(x-z)|^2}$, which is $\simeq\frac{\pi}{|(x-w)\times(x-z)|}$ when $w$ lies on the segment from $x$ to $z$: a $1/{\rm distance}$ singularity across a line, not integrable. On the segment $x-w=t\,e$, $z-w=-(1-t)e$, with $e=x-z$. Let $Q\equiv\frac{[(x'-w')\cdot e][(y-z)\cdot e]}{e^2}$, and drop the common $\frac{1}{(x'-w')^2(y-z)^2}$. The four $\mathcal A_0$ lines (3 in part 3a, 1 in part 3b) become
\begin{equation*}
\underbrace{-\frac{t}{1-t}Q-\frac{1}{1-t}Q+\frac{t}{1-t}Q}_{\rm 3a}\;\underbrace{+\frac{1}{1-t}Q}_{\rm 3b}=0.
\end{equation*}
So their sum vanishes linearly on the segment and $\mathcal A_0\times$(sum) is bounded. Part 3a or part 3b alone is not: they must stay together.

\paragraph{Group I Row II, $R^i$.} Limits, checked numerically to $10^{-6}$:
\begin{itemize}
\item $p\to\infty$: $R^i\to k^i[4\log a-\log(1-a)]$, giving $\int\frac{d^2p}{p^2}e^{ip\cdot(y-x)}\sim-2\pi\log|x-y|$ (integrable).
\item $k\to\infty$: $R^i\to2k^i\log\frac{k^2}{p^2}$, giving $\frac{(w-y)^i}{(w-y)^2}\log|w-y|$ (integrable).
\item $x,y\to w$ together: the transform is homogeneous of degree $-1$, against $r^3dr$ (integrable).
\item $p\to0$ and $p\to k$: the $\log^2$ infrared divergences of F5.
\end{itemize}

\section*{Reproducing}
In \texttt{nlo\_gluon\_paper/uv\_check/}: \texttt{uvscan2.py} (all pairs, equations and rows), \texttt{uvtarget.py} (rows with $p^+$ integrals), \texttt{intcheck.py} (Row V $d^3N_{2b\,\rm finite}$ with the integrals done), \texttt{surfscan.py} (singular curves), \texttt{make\_report.py} (this file).
\end{document}
